% English abstract; one page, structured, Times New Roman 12 pt
\ChapterStart
\begin{latin}
\begin{center}
  {\bfseries\fontsize{14pt}{17pt}\selectfont \ThesisTitleEN\par}
  \vspace{1em}
  {\bfseries\fontsize{14pt}{17pt}\selectfont Abstract}
\end{center}
\vspace{0.8em}
\noindent\textbf{Background and Aim:}
Despite pharmacological therapy and coronary revascularization, patients who have experienced myocardial infarction (MI) remain at substantial risk of recurrent cardiovascular events. Adherence to non-pharmacological secondary prevention measures, including tobacco avoidance, adequate physical activity, a cardioprotective diet, and maintaining psychological health, is essential for reducing this residual risk. Nevertheless, maintaining all these behaviors concurrently remains challenging in routine clinical practice. This study aimed to determine adherence to non-pharmacological secondary prevention measures and identify the demographic and clinical factors associated with adherence among patients with a history of MI.

\vspace{0.45em}
\noindent\textbf{Methods:}
This analytical cross-sectional study included 150 patients aged 18 years or older who had experienced MI at least one month prior to evaluation and attended the cardiology clinic of Shahid Madani Hospital in Khorramabad for follow-up during the second half of the year 1404. All eligible patients were enrolled through census sampling, and data were collected through face-to-face and telephone interviews. Demographic and clinical characteristics were recorded, while tobacco use, physical activity, adherence to the Mediterranean diet, psychological status, health literacy, illness perception, and medication adherence were assessed using standardized questionnaires. Complete adherence was defined as simultaneously meeting the criteria in all four non-pharmacological domains. Data were analyzed using descriptive statistics, comparative tests, and binary logistic regression. A P value of <0.05 was considered statistically significant.

\vspace{0.45em}
\noindent\textbf{Results:}
The participants had a mean age of $(62.1 \pm 10.7)$ years, and 70\% were men. During the 30 days preceding assessment, 82\% reported no tobacco use; 69.3\% achieved an adequate level of physical activity; 60\% adhered to the Mediterranean diet; and 46\% had a normal psychological status. Only 26 patients (17.3\%) met the criteria for adherence to all non-pharmacological measures concurrently. None had previously participated in a cardiac rehabilitation program. In multivariable analyses, men and employed participants were more likely to report smoking. Male sex, employment status, and treatment modality were independently associated with adherence to physical activity recommendations. Moderate/good economic status was associated with greater odds of dietary adherence. Marriage was associated with better psychological status, whereas female sex and more negative illness perceptions were associated with psychological distress. University education was significantly associated with complete adherence.

\vspace{0.45em}
\noindent\textbf{Conclusion:}
Complete adherence to the recommended non-pharmacological secondary prevention measures was low, and the factors associated with adherence differed across individual domains. These findings highlight the need to assess each preventive behavior separately, provide targeted support for patients with lower educational status or financial constraints, prioritize physical activity promotion among women, routinely screen for psychological distress, and establish accessible referral pathways to cardiac rehabilitation. Given the cross-sectional, single-center design, the observed associations should not be interpreted as causal.

\vspace{0.8em}
\noindent\textbf{Keywords:} \KeywordsEN
Myocardial Infarction, Secondary Prevention, Treatment Adherence and Compliance, Healthy Lifestyle, Cardiac Rehabilitation
\end{latin}
\clearpage
